% GENERATED FILE. Edit canonical lessons, then run npm run generate:books.
% canonical-source-hash: fnv1a64:aa8f3b4b11bfca43
% canonical-lessons: FR-C270-tout-a-coup, FR-C270-subitement, FR-C270-brusquement, FR-C270-sur-le-coup, FR-C270-dans-la-foulee

\chapter{Tout à coup}
\label{ch:fr-tout-a-coup}
\hlchaptermodality{270}

\begin{chapteropening}
\textbf{By the end of this chapter:} \emph{I can say that something happened all of a sudden, without warning, abruptly, on the spot or straight after.}

Complete the last lesson of chapter 270: I can say that something happened all of a sudden, without warning, abruptly, on the spot or straight after.
\end{chapteropening}

\section[tout à coup]{tout à coup — all of a sudden}

\label{lesson:FR-C270-tout-a-coup}



\begin{quote}
Before the new one: say the French for at once, then the French for while.
\end{quote}

\subsection*{You'll want to know: tout à coup}
\textbf{tout à coup} — \textquotedblleft{}all of a sudden\textquotedblright{}.

\begin{quote}
Tout à coup, il a commencé à pleuvoir.
\end{quote}

\begin{quote}
\emph{All of a sudden, it started to rain.}
\end{quote}

Close to \emph{soudain}; storytellers love it. \emph{Tout d'un coup} also means 'all at once'.

\begin{cousinweb}
\emph{coup} is inherited from Late Latin \emph{colpus} \textquotedblleft{}a blow\textquotedblright{}, from Latin \emph{colaphus} \textquotedblleft{}a punch\textquotedblright{} (EtymDB), itself borrowed from Greek.
\end{cousinweb}

\begin{grammarlens}[title={What you've built}]
One more word for telling, explaining and describing at B1.
\end{grammarlens}

\subsection*{Your turn}
\begin{itemize}
\raggedright
  \item \emph{Say it:} \emph{tout à coup}
  \item \emph{Say it:} \emph{tout à coup}, once more
  \item \emph{Say it:} a sentence of your own with \emph{tout à coup}
  \item \emph{Recall:} say the French for at once, then the French for while, then say \emph{tout à coup} again
\end{itemize}

\begin{tcolorbox}[breakable,colback=teal!4,colframe=teal!35!black,title={Before you move on}]
What does tout à coup mean? (\textquotedblleft{}All of a sudden\textquotedblright{}.) Say it once more.
\end{tcolorbox}
\section[subitement]{subitement — suddenly, without warning}

\label{lesson:FR-C270-subitement}



\begin{quote}
Before the new one: say the French for while, then the French for all of a sudden.
\end{quote}

\subsection*{You'll want to know: subitement}
\textbf{subitement} — \textquotedblleft{}suddenly, without warning\textquotedblright{}.

\begin{quote}
Le vent a subitement changé.
\end{quote}

\begin{quote}
\emph{The wind suddenly changed.}
\end{quote}

Built on \emph{subit}, 'sudden': \emph{une mort subite}.

\begin{cousinweb}
\emph{subitement} is built on \emph{subit} \textquotedblleft{}sudden\textquotedblright{}, a learned borrowing from Latin \emph{subitus} \textquotedblleft{}sudden\textquotedblright{}, from \emph{subīre} \textquotedblleft{}to come up, to steal upon\textquotedblright{} (Lewis \& Short).
\end{cousinweb}

\begin{grammarlens}[title={What you've built}]
One more word for telling, explaining and describing at B1.
\end{grammarlens}

\subsection*{Your turn}
\begin{itemize}
\raggedright
  \item \emph{Say it:} \emph{subitement}
  \item \emph{Say it:} \emph{subitement}, once more
  \item \emph{Say it:} a sentence of your own with \emph{subitement}
  \item \emph{Recall:} say the French for while, then the French for all of a sudden, then say \emph{subitement} again
\end{itemize}

\begin{tcolorbox}[breakable,colback=teal!4,colframe=teal!35!black,title={Before you move on}]
What does subitement mean? (\textquotedblleft{}Suddenly, without warning\textquotedblright{}.) Say it once more.
\end{tcolorbox}
\section[brusquement]{brusquement — abruptly, sharply}

\label{lesson:FR-C270-brusquement}



\begin{quote}
Before the new one: say the French for all of a sudden, then the French for suddenly.
\end{quote}

\subsection*{You'll want to know: brusquement}
\textbf{brusquement} — \textquotedblleft{}abruptly, sharply\textquotedblright{}.

\begin{quote}
La voiture s'est arrêtée brusquement.
\end{quote}

\begin{quote}
\emph{The car stopped abruptly.}
\end{quote}

\emph{Brusque} also describes a manner: \emph{un ton brusque}, a curt tone.

\begin{cousinweb}
\emph{brusquement} is built on \emph{brusque}, which was borrowed from Italian \emph{brusco} \textquotedblleft{}rough, sharp, sour\textquotedblright{} (EtymDB).
\end{cousinweb}

\begin{grammarlens}[title={What you've built}]
One more word for telling, explaining and describing at B1.
\end{grammarlens}

\subsection*{Your turn}
\begin{itemize}
\raggedright
  \item \emph{Say it:} \emph{brusquement}
  \item \emph{Say it:} \emph{brusquement}, once more
  \item \emph{Say it:} a sentence of your own with \emph{brusquement}
  \item \emph{Recall:} say the French for all of a sudden, then the French for suddenly, then say \emph{brusquement} again
\end{itemize}

\begin{tcolorbox}[breakable,colback=teal!4,colframe=teal!35!black,title={Before you move on}]
What does brusquement mean? (\textquotedblleft{}Abruptly, sharply\textquotedblright{}.) Say it once more.
\end{tcolorbox}
\section[sur le coup]{sur le coup — at the time, on the spot}

\label{lesson:FR-C270-sur-le-coup}



\begin{quote}
Before the new one: say the French for suddenly, then the French for abruptly.
\end{quote}

\subsection*{You'll want to know: sur le coup}
\textbf{sur le coup} — \textquotedblleft{}at the time, on the spot\textquotedblright{}.

\begin{quote}
Sur le coup, je n'ai rien compris.
\end{quote}

\begin{quote}
\emph{At the time, I did not understand anything.}
\end{quote}

The feeling of the very moment; later you see things differently.

\begin{cousinweb}
\emph{coup} is inherited from Late Latin \emph{colpus} \textquotedblleft{}a blow\textquotedblright{}, from Latin \emph{colaphus} \textquotedblleft{}a punch\textquotedblright{} (EtymDB): at the moment of the blow.
\end{cousinweb}

\begin{grammarlens}[title={What you've built}]
One more word for telling, explaining and describing at B1.
\end{grammarlens}

\subsection*{Your turn}
\begin{itemize}
\raggedright
  \item \emph{Say it:} \emph{sur le coup}
  \item \emph{Say it:} \emph{sur le coup}, once more
  \item \emph{Say it:} a sentence of your own with \emph{sur le coup}
  \item \emph{Recall:} say the French for suddenly, then the French for abruptly, then say \emph{sur le coup} again
\end{itemize}

\begin{tcolorbox}[breakable,colback=teal!4,colframe=teal!35!black,title={Before you move on}]
What does sur le coup mean? (\textquotedblleft{}At the time, on the spot\textquotedblright{}.) Say it once more.
\end{tcolorbox}
\section[dans la foulée]{dans la foulée — straight after, in the same stride}

\label{lesson:FR-C270-dans-la-foulee}



\begin{quote}
Before the new one: say the French for abruptly, then the French for at the time.
\end{quote}

\subsection*{You'll want to know: dans la foulée}
\textbf{dans la foulée} — \textquotedblleft{}straight after, in the same stride\textquotedblright{}.

\begin{quote}
Il a quitté l'école et il a trouvé un travail dans la foulée.
\end{quote}

\begin{quote}
\emph{He left school and found a job straight after.}
\end{quote}

Neutral and frequent in the press. \emph{La foulée} is a runner's stride.

\begin{cousinweb}
\emph{foulée} is built on \emph{fouler} \textquotedblleft{}to tread, to press\textquotedblright{}, which is inherited through Old French \emph{foler} from a verb built on Latin \emph{fullō} \textquotedblleft{}a fuller, a cloth-treader\textquotedblright{} (EtymDB).
\end{cousinweb}

\begin{grammarlens}[title={What you've built}]
One more word for telling, explaining and describing at B1.
\end{grammarlens}

\subsection*{Your turn}
\begin{itemize}
\raggedright
  \item \emph{Say it:} \emph{dans la foulée}
  \item \emph{Say it:} \emph{dans la foulée}, once more
  \item \emph{Say it:} a sentence of your own with \emph{dans la foulée}
  \item \emph{Recall:} say the French for abruptly, then the French for at the time, then say \emph{dans la foulée} again
\end{itemize}

\begin{tcolorbox}[breakable,colback=teal!4,colframe=teal!35!black,title={Before you move on}]
What does dans la foulée mean? (\textquotedblleft{}Straight after, in the same stride\textquotedblright{}.) Say it once more.
\end{tcolorbox}
